\subsection{问题一的求解}

\subsubsection{数据预处理}

首先从附件1读取车辆成本参数（\cref{tab:vehicle-costs}），从附件2读取排货信息。附件2包含2018年8月24日至11月23日期间从南通仓库出发的732条排货记录，按运输单号聚合为294个运输趟次。通过8层地址解析策略（特殊地址映射、正则匹配"XX市"、子串匹配城市坐标库等），从原始排货地址中提取标准城市名，共覆盖43个城市。将城市名匹配至包含155个中国城市经纬度的坐标库，计算各趟次从南通出发的单程道路距离。

\begin{table}[!h]
    \centering
    \begin{tabular}{lccc}
        \toprule[1.5pt]
        成本项目 & 4.2M & 7.6M & 9.6M \\
        \midrule[1pt]
        购置价（元）& 193{,}453 & 263{,}991 & 339{,}359 \\
        报废补贴（元）& 12{,}000 & 12{,}000 & 12{,}000 \\
        油耗（L/100km）& 13 & 17 & 20 \\
        油费（元/km）& 1.00 & 1.40 & 1.70 \\
        路桥费（元/km）& 0.70 & 1.00 & 1.30 \\
        人工（元/年）& 168{,}000 & 168{,}000 & 168{,}000 \\
        保险（元/年）& 11{,}504 & 16{,}048 & 16{,}286 \\
        审验（元/年）& 1{,}650 & 1{,}749 & 1{,}749 \\
        违章（元/年）& 60{,}833 & 60{,}833 & 60{,}833 \\
        轮胎（元/年）& 1{,}800 & 3{,}000 & 3{,}000 \\
        维保（元/年）& 84{,}000 & 96{,}000 & 96{,}000 \\
        冷藏车容量（托）& 5 & 11 & 14 \\
        \bottomrule[1.5pt]
    \end{tabular}\caption{各型号车辆成本参数（附件1）}
    \label{tab:vehicle-costs}

\end{table}

\subsubsection{费用计算结果}

基于该成本模型计算，运输总费用为\textbf{3,846,906.89元}。其中可变成本（油费+路桥费，往返）1,848,938.21元（占48.1\%），固定成本（人工+摊销+折旧，往返）1,997,968.69元（占51.9\%）。固定成本与可变成本占比接近，说明两者在运输总成本中同等重要，如\cref{fig:cost_pie}：
\begin{figure}[!h]
    \centering
    % 第一张图
    \begin{subfigure}[b]{0.26\textwidth}
        \centering
        \includegraphics[width=\textwidth]{图1_问题一_成本构成饼图.png}
        \caption{运输成本总体构成}
        \label{fig:cost_pie}
    \end{subfigure}
    \hfill % 在子图之间插入弹性间距，使它们均匀分布
    % 第二张图
    \begin{subfigure}[b]{0.36\textwidth}
        \centering
        \includegraphics[width=\textwidth]{图7_固定成本构成饼图.png}
        \caption{固定成本构成}
        \label{fig:fixed_cost}
    \end{subfigure}
    \hfill
    % 第三张图
    \begin{subfigure}[b]{0.36\textwidth}
        \centering
        \includegraphics[width=\textwidth]{图8_可变成本构成饼图.png}
        \caption{可变成本构成}
        \label{fig:variable_cost}
    \end{subfigure}
    
    \caption{运输成本构成}
    \label{fig:total_figure}
\end{figure}

按车型统计（\cref{tab:by-vehicle}），7.6M冷藏车承担了274趟运输（93.2\%），9.6M承担了20趟（6.8\%），4.2M因容量较小未被选用。

\begin{table}[!h]
    \centering
    \begin{tabular}{lcccc}
        \toprule[1.5pt]
        车型 & 趟次数 & 占比 & 总成本（元）& 平均成本（元/趟）\\
        \midrule[1pt]
        7.6M & 274 & 93.2\% & 3{,}540{,}382.67 & 12{,}921.11 \\
        9.6M & 20 & 6.8\% & 306{,}524.23 & 15{,}326.21 \\
        \midrule[1pt]
        合计 & 294 & 100\% & 3{,}846{,}906.89 & 13{,}084.72 \\
        \bottomrule[1.5pt]
    \end{tabular}\caption{按车型成本统计}
    \label{tab:by-vehicle}

\end{table}

运输距离统计（\cref{tab:distance}）显示，最短单程为南通至无锡（97.5\,km），最长为南通至乌鲁木齐（5{,}080.8\,km），平均距离1{,}254.4\,km。运距离散程度较大，反映出全国配送网络覆盖范围广泛。

\begin{table}[!h]
    \centering
    \begin{tabular}{lc}
        \toprule[1.5pt]
        统计项 & 数值 \\
        \midrule[1pt]
        最短单程道路距离 & 97.5 km（南通→无锡市）\\
        最长单程道路距离 & 5{,}080.8 km（南通→乌鲁木齐市）\\
        平均单程道路距离 & 1{,}254.4 km \\
        \bottomrule[1.5pt]
    \end{tabular}\caption{运输距离统计}
    \label{tab:distance}

\end{table}

按区域统计（\cref{tab:by-region}），东部地区以45.9\%的成本占比承载了最多的运输量，西部地区虽仅有42趟次，但平均距离达2{,}898\,km，导致单位成本远高于东部。

\begin{table}[!h]
    \centering
    \begin{tabular}{lccccc}
        \toprule[1.5pt]
        区域 & 绕行系数 & 趟次数 & 总成本（元）& 成本占比 & 平均距离（km）\\
        \midrule[1pt]
        东部 & 1.25 & 181 & 1{,}765{,}785.00 & 45.9\% & 776 \\
        中部 & 1.40 & 71 & 1{,}006{,}127.33 & 26.2\% & 1{,}501 \\
        西部 & 1.60 & 42 & 1{,}074{,}994.57 & 27.9\% & 2{,}898 \\
        \bottomrule[1.5pt]
    \end{tabular}\caption{按区域运输成本统计}
    \label{tab:by-region}

\end{table}

按城市统计的前15位（\cref{fig:cost_by_city}），乌鲁木齐以11趟次总计504{,}781.30元居首位，北京（33趟次、430{,}737.88元）和天津（22趟次、313{,}818.16元）紧随其后。远距离城市（如乌鲁木齐、哈尔滨、西安）虽趟次数较少，但单趟成本高，对总成本的贡献显著。

\begin{figure}[!h]
    \centering
    \includegraphics[width=0.8\textwidth]{图2_问题一_各城市运输成本柱状图.png}
    \caption{按城市统计的运输成本}
    \label{fig:cost_by_city}
\end{figure}

按运输时效统计（\cref{tab:by-time}），2天时效的趟次最多（113趟），1天时效平均成本最低（5{,}799.09元/趟），4天以上时效的单趟平均成本超过3万元，反映出远距离高时效运输的高成本特征。

\begin{table}[!h]
    \centering
    \begin{tabular}{lccc}
        \toprule[1.5pt]
        运输天数 & 趟次数 & 总成本（元）& 平均成本（元/趟）\\
        \midrule[1pt]
        1天 & 108 & 626{,}301.63 & 5{,}799.09 \\
        2天 & 113 & 1{,}549{,}215.13 & 13{,}709.87 \\
        3天 & 42 & 710{,}213.20 & 16{,}909.84 \\
        4天 & 25 & 784{,}451.57 & 31{,}378.06 \\
        5天 & 6 & 176{,}725.35 & 29{,}454.23 \\
        \bottomrule[1.5pt]
    \end{tabular}\caption{按运输时效成本统计}
    \label{tab:by-time}

\end{table}

\subsubsection{敏感性分析}

为识别对总成本影响最大的关键参数，基于成本模型的数学结构进行量级估算。估算逻辑为：参数变动 $\to$ 中间变量（距离、日成本等）变动比例 $\to$ 乘以该变量在总成本中的占比，得到总成本的近似变动幅度。以下数值均为该方法的粗略估计，仅用于比较各参数的相对敏感度排序。

\textbf{区域绕行系数 $\beta_r$：}绕行系数是距离计算的基础乘子（$d_r = \beta_r d_s$），直接联动可变成本（线性正比）和固定成本（通过运输天数阶梯近似线性）。若三区域绕行系数同步变动 $\pm 0.05$（按全国加权均值约 $1.38$，变动幅度约 $3.6\%$），总成本将同向变动约 $3\%$–$4\%$（$3.6\% \times$ 近1的传递系数）。西部地区（$\beta_W = 1.60$）的单趟平均成本已是东部的2.6倍，是中部的1.8倍（根据\cref{tab:by-region}计算），西部城市对 $\beta_W$ 的敏感度显著更高——若 $\beta_W$ 单独上调0.05，西部总成本增量约占全国总增量的50\%左右（假定东中部的绕行系数提升比例相同，按照比例粗略估算：$\dfrac{1}{1+\frac{1}{2.6}+\frac{1}{1.8}}\approx 51\%$）。绕行系数的合理设定是成本核算精度的首要影响因素。

\textbf{空载油耗系数 $\alpha$：}$\alpha$ 仅作用于回程油费，不影响路桥费和固定成本。由程序计算可知，回程油费为 $470{,}669$ 元，占可变成本的 $25.5\%$（$470{,}669 / 1{,}848{,}938 \approx 25.5\%$），占总成本的 $12.2\%$。$\alpha$ 每变动 $0.05$（如 $0.75 \to 0.70$），回程油费同比变动 $0.05/0.75 \approx 6.7\%$，总成本约变动 $12.2\% \times 6.7\% \approx 0.8\%$。$\alpha$ 对总量敏感度有限，但在远距离运输中累积效应显著——以乌鲁木齐为例（单程 $5{,}081\,\mathrm{km}$，7.6M车型），单趟回程油费约为 $1.40 \times 0.75 \times 5{,}081 \approx 5{,}335$ 元，$\alpha$ 从 $0.75$ 降至 $0.70$ 可单趟节约约 $356$ 元，11趟累计约 $3{,}900$ 元。

\textbf{工作天与日历天的分摊选择：}人工和轮胎按280工作天摊销，保险和审验按365日历天摊销。人工成本在固定成本中占比最大（$39.1\%$，由Python程序计算），以7.6M车型为例，人工日成本为 $168{,}000 / 280 = 600$ 元/天，若改为按 $300$ 天摊销则降至 $560$ 元/天（降幅 $6.7\%$），乘以人工在固定成本中的占比及固定成本在总成本中的占比（$51.9\%$），总成本约下降 $51.9\% \times 39.1\% \times 6.7\% \approx 1.4\%$。折旧按 $4$ 年（$4 \times 280$ 天）摊销，折旧在固定成本中占比 $25.3\%$（由Python程序计算），折旧年限从 $4$ 年延长至 $5$ 年将使日折旧降低 $20\%$，总成本影响约 $51.9\% \times 25.3\% \times 20\% \approx 2.6\%$，折旧参数变动的影响不容忽视。

\textbf{车型选择决策：}当前数据中4.2M车型（容量5托）容量过小而从未被选用，7.6M（11托）承担了93.2\%的运输量。若存在大量小批量订单（<5托），启用4.2M车型可将单趟可变成本降低约30\%–40\%。车型选择对总成本的影响取决于订单托盘分布的离散程度——托盘数越分散，车型匹配优化的空间越大。

综上，绕行系数 $\beta_r$ 是成本模型中最敏感的参数，其次为工作天摊销基准和空载系数 $\alpha$，车型选择在订单结构变化时可能成为重要因素。

\subsection{问题二的求解}

\subsubsection{优化总体效果}

在问题一成本模型基础上，对294趟运输执行CW拼车优化。使用最优参数组合（方向偏离角 $\theta_d = 90^\circ$、城市间距离上限 $d_{\max} = 500\,\mathrm{km}$、组内跨度 $\theta_g = 120^\circ$、启用跨周合并），优化结果如\cref{fig:comparison}所示。

\begin{figure}[!h]
    \centering
    \includegraphics[width=0.5\textwidth]{图3_问题二_优化前后成本对比.png}
    \caption{优化前后对比}
    \label{fig:comparison}
\end{figure}

优化后成本构成中，可变成本占48.2\%，固定成本占51.8\%，与优化前的比例基本一致，表明CW算法同时削减了可变与固定两部分成本。

\subsubsection{参数网格搜索}

为确定最优参数，对4个关键参数进行全因子网格搜索（72组组合），部分结果如\cref{tab:grid-results}所示。所有包含 $d_{\max} = 500\,\mathrm{km}$ 的参数组合均达到相同的最高节约率（55.63\%），说明城市间距离上限是决定拼车深度的关键约束。跨周合并带来了约0.7个百分点的额外节约（54.94\% vs 55.63\%），证实了跨周信息整合的价值。

\begin{table}[!h]
    \centering
    \begin{tabular}{ccccc}
        \toprule[1.5pt]
        $\theta_d$ & $d_{\max}$（km）& $\theta_g$ & 跨周合并 & 节约率 \\
        \midrule[1pt]
        90$^\circ$ & 500 & 120$^\circ$ & 是 & 55.63\% \\
        120$^\circ$ & 500 & 120$^\circ$ & 是 & 55.63\% \\
        150$^\circ$ & 500 & 120$^\circ$ & 是 & 55.62\% \\
        90$^\circ$ & 500 & 90$^\circ$ & 是 & 55.63\% \\
        90$^\circ$ & 400 & 120$^\circ$ & 是 & 55.28\% \\
        90$^\circ$ & 300 & 120$^\circ$ & 是 & 54.90\% \\
        90$^\circ$ & 500 & 120$^\circ$ & 否 & 54.94\% \\
        90$^\circ$ & 200 & 120$^\circ$ & 是 & 49.89\% \\
        \bottomrule[1.5pt]
    \end{tabular}\caption{参数网格搜索结果（代表性组合）}
    \label{tab:grid-results}

\end{table}

\subsubsection{按周优化结果}

数据跨10个自然周（周36至周46），优化后部分自然周的路线数和成本如\cref{tab:week-results}所示。周41覆盖34个城市、路线最多（31条）、成本最高（428{,}528.47元），与该周订单集中度较高有关。

\begin{table}[!h]
    \centering
    \begin{tabular}{lccc}
        \toprule[1.5pt]
        周次 & 路线数 & 总成本（元）& 覆盖城市数 \\
        \midrule[1pt]
        周36 & 9 & 99{,}009.20 & 11 \\
        周37 & 5 & 126{,}887.30 & 6 \\
        周41 & 31 & 428{,}528.47 & 34 \\
        周42 & 18 & 221{,}296.37 & 22 \\
        周43 & 9 & 112{,}435.70 & 9 \\
        周44 & 20 & 264{,}463.84 & 26 \\
        周45 & 16 & 259{,}481.53 & 25 \\
        周46 & 19 & 194{,}792.47 & 20 \\
        \bottomrule[1.5pt]
    \end{tabular}\caption{各周拼车优化结果}
    \label{tab:week-results}

\end{table}

\subsubsection{典型拼车路线}

\cref{tab:example-routes}列举了几条典型的拼车路线。浙中南线将宁波、温岭、丽水、南昌、衢州、金华、义乌等7个城市合并为一条路线，9.6M车型装载13.8托，运输5天，总成本14{,}439.94元（若分别运输，仅宁波至南通的往返可变成本即超过3{,}000元）。东北线（沈阳—长春—吉林）使用9.6M车型装载11.5托，7天往返成本22{,}131.00元。

\begin{table}[!h]
    \centering
    \begin{tabular}{lcccl}
        \toprule[1.5pt]
        路线 & 车型 & 托数 & 天数 & 成本（元）\\
        \midrule[1pt]
        宁波→温岭→丽水→南昌→衢州→金华→义乌（浙中南线）& 9.6M & 13.8 & 5 & 14{,}439.94 \\
        南京→义乌→金华→丽水→温州→临海→宁波（浙东线）& 9.6M & 13.1 & 4 & 10{,}909.39 \\
        沈阳→长春→吉林（东北线）& 9.6M & 11.5 & 7 & 22{,}131.00 \\
        大连→盘锦→沈阳（辽宁线）& 9.6M & 11.6 & 5 & 16{,}015.49 \\
        太原→榆林→呼和浩特（晋蒙线）& 7.6M & 6.8 & 8 & 21{,}965.27 \\
        潍坊→聊城→郑州（鲁豫线）& 7.6M & 8.7 & 5 & 13{,}192.18 \\
        \bottomrule[1.5pt]
    \end{tabular}\caption{典型拼车路线示例}
    \label{tab:example-routes}

\end{table}

\subsubsection{敏感性分析}

\cref{tab:grid-results}已呈现了部分参数组合的网格搜索结果（由Python程序实际计算得到），以下从参数敏感性角度进行归纳。

\textbf{城市间距离上限 $d_{\max}$（第一敏感参数）：}$d_{\max}$ 直接决定了可合并城市对的候选集大小。当 $d_{\max}$ 从200\,km增至300\,km时，节约率从49.9\%跃升至54.9\%（提升5.0个百分点）；从300\,km增至400\,km时升至55.5\%（仅提升0.6个百分点）；从400\,km增至500\,km时升至55.63\%（仅提升0.15个百分点）。呈现典型的边际递减特征——300\,km以内的放宽释放了大量省内短途合并机会，而400\,km以上的进一步放宽收效甚微。在实际运营决策中，$d_{\max}$ 不仅影响成本节约，还关系到拼车路线的驾驶时长和驾驶员工作强度，需在节约效果与运营可行性之间权衡。

\textbf{方向偏离角 $\theta_d$ 与组内跨度 $\theta_g$（次敏感参数）：}在 $d_{\max} = 500\,\mathrm{km}$ 的宽松距离约束下，$\theta_d$ 从90°放宽至150°、$\theta_g$ 从90°放宽至120°，节约率几乎无变化（均稳定在55.62\%–55.63\%）。这说明当距离约束足够宽松时，方向约束已不再构成瓶颈——满足距离约束的城市对通常也满足方向约束。然而，在 $d_{\max}$ 较紧（如200\,km）时，$\theta_d = 90^\circ$ 与 $\theta_d = 150^\circ$ 之间存在可见差异（差距约1–2个百分点），方向约束在短途合并中具有实质筛选作用。

\textbf{跨周合并（边际增益参数）：}启用跨周合并带来约0.7个百分点的额外节约（54.94\% $\to$ 55.63\%），增幅虽小但绝对值可观（对应约2.6万元）。跨周合并的价值取决于订单的时间分布——若一周内某方向仅有极少数订单，跨周信息整合可将其与相邻周的顺路订单合并，避免单独发车。

\textbf{容量约束 $Q_v$ 的隐性影响：}7.6M车型容量为11托，9.6M为14托。在网格搜索最优解中，约68\%的合并路线选用9.6M车型（因其承载能力更强、可合并更多订单），说明容量约束是限制拼车深度的结构性因素。若车队中9.6M车型占比低于最优解中的使用比例，实际节约率将低于理论值。

综上，$d_{\max}$ 是拼车优化的第一敏感参数且呈边际递减特征，$\theta_d$ 和 $\theta_g$ 仅在距离约束收紧时发挥作用，跨周合并和容量约束分别从时间和空间维度影响优化上限。

\subsection{问题三的求解}

\subsubsection{数据规模与仓库分类}

附件3包含536{,}208条运单记录，附件4包含222{,}378条派车记录，覆盖152个城市、1{,}186辆车牌和88个仓库编码。依据有无派车记录，将88个仓库分为三类：A类仓库54个（有运单且有派车记录），是评估和优化的核心对象；B类仓库16个（仅有派车记录无运单）；C类仓库18个（仅有运单无派车记录）。最终评估覆盖54个仓库，总车辆日数为3{,}695。

\subsubsection{六维指标得分}

54个仓库在六个指标上的绝对得分描述性统计如\cref{tab:desc}所示。时效达标均分最高（98.6分），且标准差最小（4.7），说明各仓库在时效方面普遍表现良好。装载饱和的均值最低（37.7分），反映出多数仓库的车辆装载量远低于行业p75水平、运力浪费较为普遍。冷链合规的标准差最大（39.5），是仓库间差异最大的维度。线路集中均分仅为29.1分，标准差27.9，反映出多数仓库的线路结构较为分散、需求碎片化程度高，缺乏成熟的骨干线路网络。

\begin{table}[!h]
    \centering
    \begin{tabular}{lcccc}
        \toprule[1.5pt]
        指标 & 均值 & 标准差 & 最小值 & 最大值 \\
        \midrule[1pt]
        配送密度 & 56.6 & 24.6 & 8.3 & 98.7 \\
        路线潜力 & 65.3 & 16.3 & 50.0 & 100.0 \\
        时效达标 & 98.6 & 4.7 & 71.8 & 100.0 \\
        线路集中 & 29.1 & 27.9 & 1.7 & 100.0 \\
        冷链合规 & 59.5 & 39.5 & 0.0 & 100.0 \\
        装载饱和 & 37.7 & 25.0 & 0.0 & 84.9 \\
        \bottomrule[1.5pt]
    \end{tabular}\caption{各指标绝对得分描述性统计}
    \label{tab:desc}

\end{table}

\subsubsection{熵权法权重计算}

基于54个仓库在6项指标上的绝对得分矩阵，运用熵权法计算客观权重，结果如\cref{tab:entropy}所示。路线合并潜力获得了最高权重（34.91\%），其信息熵仅为0.8397——意味着各仓库在路线规划上存在显著差异，该指标对评价的区分能力最强。线路集中权重居第二（23.17\%），信息熵为0.8936，反映出各仓库在需求结构的集中度上差异同样突出。装载饱和与冷链合规的权重相近（分别为15.94\%和15.83\%），同属区分能力较强的指标。相反，时效达标的信息熵高达0.9937，权重仅1.37\%——几乎所有仓库的送达时效都接近100\%，该指标无法区分仓库的优劣。

\begin{table}[!h]
    \centering
    \begin{tabular}{lccc}
        \toprule[1.5pt]
        指标 & 信息熵 $e_j$ & 差异系数 $d_j$ & 权重 $w_j$ \\
        \midrule[1pt]
        配送密度 & 0.9597 & 0.0403 & 8.76\% \\
        路线合并潜力 & 0.8397 & 0.1603 & 34.91\% \\
        时效达标率 & 0.9937 & 0.0063 & 1.37\% \\
        线路集中度 & 0.8936 & 0.1064 & 23.17\% \\
        冷链合规率 & 0.9273 & 0.0727 & 15.83\% \\
        装载饱和度 & 0.9268 & 0.0732 & 15.94\% \\
        \bottomrule[1.5pt]
    \end{tabular}\caption{熵权法权重计算结果}
    \label{tab:entropy}

\end{table}

为直观表示各指标对系统得分的贡献，\cref{fig:sys-score}直观展示了各指标的权重占比。
\begin{figure}[!h]
    \centering
    \includegraphics[width=0.5\textwidth]{图4_问题三_熵权法权重饼图.png}
    \caption{各指标权重占比}
    \label{fig:sys-score}
\end{figure}

\subsubsection{系统综合得分}

系统加权综合得分为\textbf{50.0分（满分为100分）}，评级为"需改进"。各指标对系统得分的贡献如\cref{tab:sys-score}所示。路线合并潜力贡献了23.6分（占总分的47.2\%），是决定系统得分的最主要因素。装载饱和平均仅48.1分、冷链合规平均仅49.7分，两者合计贡献15.6分，是拖累整体评分的关键短板。

\begin{table}[!h]
    \centering
    \begin{tabular}{lccc}
        \toprule[1.5pt]
        指标 & 权重 & 加权均分 & 贡献值 \\
        \midrule[1pt]
        配送密度 & 8.76\% & 65.8 & 5.8 \\
        路线合并潜力 & 34.91\% & 67.6 & 23.6 \\
        时效达标率 & 1.37\% & 99.0 & 1.4 \\
        线路集中度 & 23.17\% & 16.2 & 3.8 \\
        冷链合规率 & 15.83\% & 49.7 & 7.9 \\
        装载饱和度 & 15.94\% & 48.1 & 7.7 \\
        \midrule[1pt]
        合计 & 100.00\% & — & 50.0 \\
        \bottomrule[1.5pt]
    \end{tabular}\caption{系统综合得分构成}
    \label{tab:sys-score}

\end{table}

\subsubsection{仓库排名分析}

全部54个仓库的合理性排名如\cref{fig:wh-ranking}所示。SH6仓库以84.4分位居第一，各项指标表现突出（配送密度75.0、路线潜力88.4、线路集中100.0、冷链合规100.0）；ZZ2仓库以26.8分垫底，各维度均存在不足（配送密度47.2、路线潜力50.0、冷链合规0.0、装载饱和0.0）。

TOPSIS辅助排名与综合得分排名的Spearman秩相关系数为0.9482（$p = 0.0000$），表明两种评估方法给出的排名高度一致，熵权法的评估结论是稳健可信的。

\begin{figure}[!h]
    \centering
    \includegraphics[width=0.8\textwidth]{图5_问题三_仓库综合得分排名.png}
    \caption{仓库综合得分排名（前15名）}
    \label{fig:wh-ranking}
\end{figure}

\subsubsection{冷链合规系统性风险}

23个仓库（占42.6\%）的冷链合规得分低于50分，存在冷链药品被违规使用常温车配送的风险。其中N02仓库（625单冷链订单，合规率0\%）、NJ1（41单冷链，0\%）、ZZ2（39单冷链，0\%）、HM1（102单冷链，0\%）和FJ4（1单冷链，0\%）五个仓库的冷链合规率均为0\%，属于严重不合规。SX1（428单冷链，合规率21.3\%）和WH1（574单冷链，合规率44.3\%）等大型仓库也存在大规模冷链药品使用常温车的问题，亟需系统性整改。

\subsubsection{三层递进优化}

对54个仓库的派车数据执行三层递进优化，结果如\cref{fig:opt}所示。第1层TSP路径排序实现5.6\%的节约，第2层和第3层的合并与车型匹配进一步将节约率推升至36.8\%。优化后路线数从3{,}379个车辆日降至1{,}451条合并路线，车辆日减少57.1\%，总成本从1{,}135.3万元降至717.5万元，节约417.9万元。

\begin{figure}[!h]
    \centering
    \includegraphics[width=0.6\textwidth]{图6_问题三_三层优化瀑布图.png}
    \caption{三层递进优化效果}
    \label{fig:opt}
\end{figure}
\subsubsection{敏感性分析}

对CW算法中城市间距离阈值参数进行敏感性分析，结果如\cref{tab:sens-dist}所示。当阈值从200\,km增至400\,km时，节约率从35.2\%升至36.8\%，增幅显著；从400\,km增至600\,km时，节约率仅小幅升至37.3\%，说明400\,km附近的阈值已能捕获主要的合并机会，继续放宽的边际收益递减。

\begin{table}[!h]
    \centering
    \begin{tabular}{lccc}
        \toprule[1.5pt]
        距离阈值（km）& 路线数 & 成本（万元）& 节约率 \\
        \midrule[1pt]
        200 & 1{,}533 & 735.3 & 35.2\% \\
        400 & 1{,}451 & 717.5 & 36.8\% \\
        600 & 1{,}425 & 711.5 & 37.3\% \\
        \bottomrule[1.5pt]
    \end{tabular}\caption{城市间距离阈值敏感性分析}
    \label{tab:sens-dist}

\end{table}

除CW距离阈值（由Python程序实际计算）外，评估体系的得分和权重对核心假设的敏感性同样值得关注。

\textbf{熵权法权重对极端仓库的敏感性（由Python程序实际计算）：}去除冷链合规得分为0的5个极端仓库（N02、NJ1、ZZ2、HM1、FJ4）后，重新对49个仓库执行熵权法计算，权重对比如\cref{tab:entropy-sens}所示。剔除零分仓库后，冷链合规指标的信息熵上升（差异系数下降），其权重从 $15.83\%$ 降至 $12.23\%$（降幅 $-3.61$ 个百分点），释放的权重被其余五指标分摊——线路集中上升 $+2.13$ 个百分点至 $25.30\%$，装载饱和上升 $+1.40$ 个百分点至 $17.34\%$，配送密度上升 $+0.77$ 个百分点至 $9.53\%$。路线合并潜力权重几乎不变（$34.91\% \to 33.96\%$），仍为最高权重指标。权重排序完全保持：路线潜力 $>$ 线路集中 $>$ 装载饱和 $>$ 冷链合规 $>$ 配送密度 $>$ 时效达标。熵权法对极端样本的稳健性优于直觉预期——极端仓库的存在并未扭曲权重排序，仅对绝对值有一定影响。

\begin{table}[!h]
    \centering
    \begin{tabular}{lccc}
        \toprule[1.5pt]
        指标 & 剔除前权重（54仓）& 剔除后权重（49仓）& 变动 \\
        \midrule[1pt]
        配送密度 & 8.76\% & 9.53\% & $+0.77$ pp \\
        路线合并潜力 & 34.91\% & 33.96\% & $-0.95$ pp \\
        时效达标率 & 1.37\% & 1.63\% & $+0.26$ pp \\
        线路集中度 & 23.17\% & 25.30\% & $+2.13$ pp \\
        冷链合规率 & 15.83\% & 12.23\% & $-3.61$ pp \\
        装载饱和度 & 15.94\% & 17.34\% & $+1.40$ pp \\
        \bottomrule[1.5pt]
    \end{tabular}\caption{熵权法极端仓库剔除前后权重对比}
    \label{tab:entropy-sens}
\end{table}

\textbf{评分函数基准值的敏感性（量级估算）：}配送密度采用中间型梯形隶属函数（市内最优区间 $[12, 18]$ 点/日，干线 $[3, 5]$ 点/日），其敏感性主要体现在最优区间边界而非线性斜率。若干线下界从3上调至3.5，部分干线型仓库的得分将显著下降；但配送密度权重仅 $8.76\%$，系统总分变动不超过 $1$ 分。线路集中度的评分完全由数据驱动（基于各仓库自身线路分布的信息熵归一化），无需外部基准假设，因此不受基准设定的影响。

\textbf{冷链合规作为硬约束的敏感性：}当前评估将冷链合规作为软指标，若改为硬约束（冷链合规率低于 $50\%$ 的仓库直接判定为不合格），则 $23$ 个仓库（$42.6\%$）将触发整改。此为政策场景分析而非参数微调：优化优先级从“降本”转变为“合规达标”，三层优化的第2层（同城同日合并、冷链常温分离）优先度大幅提升。整体节约率受冷链车辆资源约束，预计下降 $5$–$10$ 个百分点（估算依据：冷链车辆数约为常温车辆的 $60\%$，合规约束使冷链车辆的可用容量成为瓶颈，约 $30\%$ 的常温配送点可能因合并路线中的冷链需求而被迫分拆）。

综上，问题三的评估体系对极端仓库样本和评分基准的设定具有一定敏感性，但核心结论在合理参数变动范围内保持稳健：（1）路线规划能力始终是首要区分因素（权重排序不受极端样本影响）；（2）路线潜力与线路集中始终是最具区分度的前两位指标。CW距离阈值的敏感性呈边际递减特征（已由程序计算验证），冷链合规的硬约束化是影响优化空间的最大政策变量。
